\section{Dinero nuevo, dinero viejo y oportunidades para el país y la nación}

La sección anterior trató la entrada de la empresa en el sistema de las armadoras; esta sección pasa a una narrativa más amplia sobre la globalización y la industria. Li Yifan (李一帆) habla del dinero nuevo y del dinero viejo, y también de las oportunidades para el país, la nación y la industria. El emprendimiento en tecnología dura es distinto del internet de consumo: con frecuencia está ligado a la capacidad manufacturera del país, a la integridad de la cadena industrial, a la comunicación con clientes de todo el mundo y a las diferencias entre culturas de capital.

\begin{figure}[H]
\centering
\includegraphics[width=0.94\textwidth]{new-money-old-money.png}
\caption{Dinero nuevo y dinero viejo: en la comunicación global, cada cultura de capital y de industria tiene su propio lenguaje de confianza. Diagrama conceptual propio, elaborado a partir de la conversación entre 02:38:34 y 03:02:16.}
\end{figure}

\begin{knowledgebox}{Lectura de la figura: globalizarse no es solo traducir el idioma}
El dinero viejo valora la historia, el orden y la estabilidad; el dinero nuevo valora la velocidad, el crecimiento y la tecnología. Cuando una empresa china de tecnología dura se globaliza, no solo tiene que explicar su producto, sino también demostrar que es confiable, que piensa a largo plazo y que puede cumplir sus entregas.
\end{knowledgebox}

\subsection{País, nación y oportunidades industriales}

Después de la globalización, esta subsección vuelve a la narrativa industrial más amplia. Tanto al principio como al final de la conversación aparece la expresión «país, nación, oportunidad». No se trata de una consigna abstracta, sino de una realidad industrial: los vehículos de nueva energía, los sensores, los chips, la cadena de suministro, la manufactura y el mercado global determinan en conjunto si una empresa de tecnología dura tiene una oportunidad. Algunas oportunidades no las crea un emprendedor individual, sino que las otorgan la época y la industria.

\begin{figure}[H]
\centering
\includegraphics[width=0.92\textwidth]{national-industrial-opportunity.png}
\caption{País, nación y oportunidades industriales: el emprendimiento en tecnología dura suele surgir de la suma de la época, la cadena industrial y la capacidad del país. Diagrama conceptual propio, elaborado a partir de la conversación entre 00:00:04 y 00:00:59 y entre 02:38:34 y 03:02:16.}
\end{figure}

\begin{warningbox}{La gran narrativa no sustituye la capacidad de producto}
Las oportunidades para el país y la nación abren una ventana industrial, no una victoria gratuita. La empresa todavía tiene que aprovechar esa ventana con tecnología, costos, calidad, clientes y capacidad organizacional.
\end{warningbox}

\begin{knowledgebox}{El examen final de una empresa de tecnología dura}
\begin{tabularx}{\textwidth}{p{0.22\textwidth}p{0.34\textwidth}X}
\toprule
Examen & Pregunta & Criterio de aprobación \\
\midrule
Examen técnico & ¿El desempeño es realmente superior? & Pruebas objetivas, escenarios de clientes, confiabilidad a largo plazo. \\
Examen de costos & ¿Puede reducir costos de forma sostenida? & BOM, rendimiento de producción, cadena de suministro y escalamiento. \\
Examen de clientes & ¿Las armadoras están dispuestas a adoptarlo? & Certificación, nominación como proveedor, producción en serie y recompra. \\
Examen de capital & ¿Puede atravesar los ciclos? & Financiamiento, flujo de efectivo, salida a bolsa y confianza del mercado. \\
Examen de globalización & ¿Puede cumplir entregas entre culturas distintas? & Clientes en el extranjero, cumplimiento normativo, marca y red de servicio. \\
\bottomrule
\end{tabularx}
\end{knowledgebox}

\begin{knowledgebox}{Lista de verificación para la comunicación global}
\begin{tabularx}{\textwidth}{p{0.22\textwidth}p{0.34\textwidth}X}
\toprule
Dimensión & Qué hay que demostrar & Materiales típicos \\
\midrule
Credibilidad técnica & Que el desempeño y la confiabilidad no son discurso publicitario & Informes de pruebas, casos de clientes, certificación de grado automotriz. \\
Credibilidad comercial & Que la empresa puede suministrar y dar servicio a largo plazo & Finanzas, capacidad de producción, red de posventa. \\
Credibilidad cultural & Que entiende las culturas de distintos clientes y capitales & Equipo local, estilo de comunicación, compromiso a largo plazo. \\
Credibilidad de cumplimiento & Que cumple las normas locales y los requisitos de la cadena de suministro & Revisión de cumplimiento normativo, procesos de datos y de seguridad. \\
\bottomrule
\end{tabularx}
\end{knowledgebox}

\subsection{Resumen de la sección}

El emprendimiento en tecnología dura es una historia de empresa y, al mismo tiempo, una historia de industria y de capacidad nacional. Para entender a Hesai (禾赛), hay que mirar a la vez los detalles técnicos y la ventana que abrió la época.
